\chapter*{Ringraziamenti}
\addcontentsline{toc}{chapter}{Ringraziamenti}

Ho iniziato questo scritto con una frase celebre del crittografo Bruce Schneier e proprio come dice quella frase, questo scritto non è un prodotto ma un processo.
Un processo che racchiude un po' tutto un percorso importante della mia vita, un percorso durato 3 anni.

Qui dentro si può intravedere un po' la persona che sono, la passione e l'impegno che ci metto nelle cose che amo fare.
Poter fare qualcosa che si ama non è per nulla scontato al giorno d'oggi e per questo vorrei ringraziare chi mi ha dato la possibilità di farlo e chi mi ha sostenuto nel farlo.

Ringrazio il prof. Marchetti che per primo ha accolto la mia idea, dandomi la possibilità di studiare, progettare e realizzare il lavoro che sta alla base di questo elaborato.
Una possibilità che però senza l'aiuto e la guida del dott. Galli non sarei stato in grado di fare. 
Grazie a lui ho imparato molte cose nel ambito del Machine Learning e Cybersecurity, ma anche quello lavorativo e di gruppo. 
Infatti sono stato accolto in un fantastico gruppo con delle bellissime persone le quali farò fatica a dimenticarmi, grazie ragazzi e grazie dottore.

Tutti loro mi hanno accompagnato in questi ultimi mesi, ma come ho detto, questo percorso è durato molto di più.
Non dimenticherò mai tutte le lezioni, i momenti e gli esami passati con il mio gruppo di amici.
Amici che se non avessi mai scelto questa strada non avrei mai incontrato e, senza saperlo, la mia vita sarebbe stata un po' più vuota.
Scusate ma per fortuna siete tanti per descrivervi e ringraziarvi tutti per le bellissime persone che siete, dovrei scrivere un altro elaborato e per ciò mi limito purtroppo a ringrazziarvi immensamente.
A loro si uniscono però anche il mio gruppo di amici e la mia migliore amica, che ormai dalle superiori mi aiutano a svagare e a staccare un po' la spina nei momenti più stressanti.
Stessa cosa vale anche per voi, io per quest'anno ho dato con la scrittura.

Fino a qui ho parlato di professori, colleghi e amici. Tutti loro sono tante piccole famiglie a cui voglio bene in mille modi diversi.
Tuttavia, anche se le definisco come tali, soltanto una è quella che mi ha sostenuto fin da quando sono nato, la mia vera famiglia.
\begin{quote}
Una famiglia del tutto particolare, ma bella.\\
Una famiglia amorevole, piena di vita e felicità.\\
Una famiglia di cui posso essere orgoglioso e fiero di farne parte.
\end{quote}
Ringrazio mio Padre che da lui ho preso un po' il mio perfezionismo, i miei Nonni e in particolare Cumpà da cui ho preso l'inventiva, i miei Zii e i miei cugini che si sono sempre interessati al mio percorso e non vedono l'ora di festeggiarmi.
Ovviamente non mi sono dimenticato di mio Fratello, colui che ancor prima di finire le superiori mi ha spronato a fare questo percoso e ne sono immensamente grato. Senza di lui non sarei mai arrivato a scrivere queste belle parole, grazie.
Ringrazio anche colei che da poco è entrata in questa piccola famiglia, che tanto piccola non è, rimanendo al mio fianco supportandomi e sopportandomi anche nei momenti difficili. Grazie tanto amore mio grande.

In fine non ringrazio la Tata Manu soltanto per la bellissima zia qual è, ma la ringrazio anche perchéé 3 anni fa, al inizio di questo cammino, mi diede questo messaggio.
Un messaggio scritto dalla persona più importante per me, colei che già dal 2012 credebbe in me e nei miei sogni. Colei che da qualche parte lassù mi starà guardando.

\begin{quote}
    \textit{!! Alessandro e' cresciuto molto di testa, di corpo, di tutto e' troppo forte quando e' a casa solo senza Marco ha una pacatezza disarmante e spero di non smentirmi fara' l'ingegnere !! il ns piccolo inventore..}
\end{quote}

Grazie Mamma.
